\section{The quadratic gauge from the inverse equation}\label{sec:gauge}

The polynomial realization theorem rests on one explicit construction.
Let
\begin{equation}
 G(S)=g_1S+g_2S^2+\cdots+g_NS^N\in K[S],
 \qquad g_1g_3g_N\neq0.
 \label{eq:seed}
\end{equation}
The construction is easiest to read backwards, from the inverse problem.
Begin with birational coordinates \((\Pi,S,Q)\) and put
\begin{equation}
 D=1-SQ+\Pi S^2.
 \label{eq:D-compact}
\end{equation}
For the prescribed one-variable seed, define
\begin{equation}
 G_\Pi(S)=g_1S+\Pi(g_2S^2+g_3S^3)
          +\sum_{k=4}^Ng_k\Pi^kS^k
 \label{eq:Gpi}
\end{equation}
and
\begin{equation}
 \beta(\Pi,S)
 =\frac{G_\Pi'(S)/g_1-1-\Pi S^2}{S}.
 \label{eq:beta}
\end{equation}
The numerator is divisible by \(S\), so \(\beta\) is a polynomial.  We now
choose the second target coordinate to be
\begin{equation}
 B=Q+\beta(\Pi,S).
 \label{eq:B-compact}
\end{equation}
This choice is the whole gauge.  Indeed,
\[
 D
 =1-S(B-\beta(\Pi,S))+\Pi S^2
 =\frac{G_\Pi'(S)}{g_1}-BS.
\]
Define the third target coordinate by
\begin{equation}
 C=\frac{2G_\Pi(S)}{g_1}-BS^2.
 \label{eq:C-compact}
\end{equation}
Then the inverse equation
\begin{equation}
 E_{\Pi,B,C}(S)
 =G_\Pi(S)-\frac{g_1}{2}(BS^2+C)=0
 \label{eq:inverse}
\end{equation}
holds by construction, and its derivative satisfies
\begin{equation}
 D=\frac{G_\Pi'(S)}{g_1}-BS
 =\frac{E_{\Pi,B,C}'(S)}{g_1}.
 \label{eq:D-derivative}
\end{equation}

This identity already gives the root-to-source reconstruction.  Fix
\((\pi,b,c)\in K^*\times K^2\), let \(s\) be a root of
\(E_{\pi,b,c}\), and suppose \(E_{\pi,b,c}'(s)\) is invertible.  Then the
source point is recovered by the single formula
\begin{equation}
\boxed{
\begin{gathered}
 d=\frac{E_{\pi,b,c}'(s)}{g_1},\qquad
 Q=b-\beta(\pi,s),\qquad t=d^{-1},\\
 x=sd^{-1},\qquad y=Q-\pi s,\qquad q=\pi d,\\
 z=d^2\left(q-\frac{g_1}{g_3}y^2(1+3t)\right).
\end{gathered}}
\label{eq:boxed-root-reconstruction}
\end{equation}
The only localization is the inverse of the derivative unit \(d\).  The
proof of Theorem~\ref{thm:localized-gauge} verifies over an arbitrary
commutative test algebra that this formula and extraction of the marked root
are mutually inverse.

The Jacobian calculation is already visible in these coordinates.  At fixed
\(\Pi\), differentiating \eqref{eq:B-compact}--\eqref{eq:C-compact} gives
\begin{equation}
 \det\frac{\partial(B,C)}{\partial(S,Q)}=-2D.
 \label{eq:plane-jacobian}
\end{equation}
Thus any birational source chart with Jacobian \(D^{-1}\) will produce a
constant Jacobian \(-2\).  What remains is to find such a chart for which
\(B\) and \(C\) pull back to polynomials.

\subsection{The polynomial pullback}

In source coordinates \((x,y,z)\), set
\begin{equation}
 t=1+xy,
 \qquad
 q=t^2z+\frac{g_1}{g_3}y^2(1+3t).
 \label{eq:tq}
\end{equation}
On \(t\neq0\), define
\begin{equation}
 \Pi=tq,\qquad S=\frac{x}{t},\qquad Q=y+xq.
 \label{eq:source-chart}
\end{equation}
Then
\[
 D=1-SQ+\Pi S^2=t^{-1},
\]
and direct differentiation gives
\begin{equation}
 \det\frac{\partial(\Pi,S,Q)}{\partial(x,y,z)}=D^{-1}.
 \label{eq:source-chart-jacobian}
\end{equation}
Consequently, the two Jacobian factors are
\[
 D^{-1}\cdot(-2D)=-2.
\]

The particular polynomial \(q\) in \eqref{eq:tq} is chosen so that the
compact coordinates \eqref{eq:B-compact}--\eqref{eq:C-compact} have no
denominators after this substitution.  Expanding their pullbacks gives the
polynomial map \(F_G=(\Pi,B,C):\A^3_K\to\A^3_K\):
\begin{align}
 \Pi={}&tq,\notag\\
 B={}&y+3\frac{g_3}{g_1}xq+2\frac{g_2}{g_1}tq
       +\sum_{k=4}^N k\frac{g_k}{g_1}
          t^2x^{k-2}q^k,\notag\\
 C={}&x(5-3t)-\frac{g_3}{g_1}x^3z
       -\sum_{k=4}^N(k-2)\frac{g_k}{g_1}(xq)^k.
 \label{eq:gauge-map}
\end{align}
The long formula is therefore the polynomial pullback of the simple inverse
problem above, not a formula guessed in the source coordinates.  Once
\eqref{eq:tq} is fixed, the pullback is forced coefficient by coefficient.
Appendix~\ref{app:gauge-identities} records that expansion in full.

\subsection{Motivation: the cubic seed and its quotient}

For \(G(S)=S+S^3\), formula \eqref{eq:gauge-map} specializes coordinatewise
to the map announced by Alp\"oge \cite{Alpoge}.  We use that map only as the
cubic seed of the construction; its determinant and collision now also have
an independent Isabelle/HOL formalization \cite{IsabelleAFP}.  The geometric
motivation is normalized multiplication of a marked linear factor and a
quadratic factor of a binary cubic: the three marked roots give the three
inverse branches, while a nonzero resultant supplies local invertibility
\cite{TaoDigest}.  Shaska's mixed-sign grading and quotient viewpoint gives
a complementary explanation of the Jacobian cancellation on the contracted
quotient \cite{Shaska}.  Neither viewpoint replaces the derivative-unit
formula \eqref{eq:boxed-root-reconstruction}, which is what identifies the
complete fiber functor for an arbitrary inverse polynomial.

Thus the three-variable inverse problem is governed by the single equation
\(E_{\Pi,B,C}(S)=0\), and the same derivative that controls reconstruction
is exactly the factor cancelled in the Jacobian.

\section{Localized fiber reconstruction and geometric degree}
\label{sec:reconstruction}

For $(\pi,b,c)\in K^*\times K^2$, write $\mathcal F_{\pi,b,c}$ for the fiber
functor
\[
 \mathcal F_{\pi,b,c}(A)
 =\{(x,y,z)\in A^3:F_G(x,y,z)=(\pi,b,c)\}
\]
on commutative $K$-algebras.

\begin{lemma}[Derivative units over test algebras]\label{lem:derivative-unit}
Let $E\in K[S]$ be separable, let $A$ be a commutative $K$-algebra, and let
$s\in A$ satisfy $E(s)=0$.  Then $E'(s)$ is a unit of $A$.
\end{lemma}

\begin{proof}
Choose $U,V\in K[S]$ with $UE+VE'=1$.  Evaluating at $s$ gives
$V(s)E'(s)=1$.  Since $A$ is commutative, this is a two-sided inverse.
\end{proof}

\begin{theorem}[General localized quadratic-gauge fiber theorem]
\label{thm:localized-gauge}
The map $F_G$ satisfies
\[
 \det DF_G=-2,
 \qquad
 \gdeg(F_G)=N.
\]
Its generic inverse equation is $E_{\Pi,B,C}(S)=0$.  Fix
$(\pi,b,c)\in K^*\times K^2$, put $E=E_{\pi,b,c}$, and let
$\overline{E'}$ be the class of $E'$ in $K[S]/(E)$.  Then
\begin{equation}
 \boxed{
 F_G^{-1}(\pi,b,c)
 \simeq
 \operatorname{Spec}\left((K[S]/(E))[1/\overline{E'}]\right).}
 \label{eq:localized-fiber}
\end{equation}
Equivalently, for every commutative $K$-algebra $A$ there is a natural
bijection
\begin{equation}
 \operatorname{Hom}_{K\text{-alg}}
 \left((K[S]/(E))[1/\overline{E'}],A\right)
 \simeq \mathcal F_{\pi,b,c}(A).
 \label{eq:localized-natural-fiber-equivalence}
\end{equation}
\end{theorem}

\begin{proof}
The construction above gives the identities
\begin{equation}
 B=Q+\beta(\Pi,S),
 \qquad
 C=\frac{2G_\Pi(S)}{g_1}-BS^2.
 \label{eq:marked-line-identities}
\end{equation}
Equations \eqref{eq:source-chart-jacobian} and
\eqref{eq:plane-jacobian} give \(\det DF_G=-2\) on the dense open set
\(t\neq0\), and hence everywhere.
For the cubic seed, this agrees with the normalized marked-factor
interpretation above.  The displayed calculation proves that the same
cancellation persists for every admissible seed \(G\).

We now identify the entire fiber.  Let $A$ be an arbitrary commutative
$K$-algebra.  An $A$-point $(x,y,z)$ above $(\pi,b,c)$ satisfies
\[
 tq=\pi.
\]
Because $\pi$ is a unit, both $t$ and $q$ are units.  The chart coordinates
\[
 S=xt^{-1},
 \qquad Q=y+xq,
 \qquad D=t^{-1}
\]
are therefore defined on every point of the fiber.  Equations
\eqref{eq:marked-line-identities} and \eqref{eq:D-derivative} give
\[
 Q+\beta(\pi,S)=b,
 \qquad E(S)=0,
 \qquad D=\frac{E'(S)}{g_1}.
\]
Thus every source point gives a root $S$ of $E$ for which $E'(S)$ is a unit.

Conversely, suppose $s\in A$ satisfies $E(s)=0$ and $E'(s)\in A^*$.  Put
the variables \(d,Q,t,x,y,q,z\) equal to the expressions in the boxed
reconstruction formula \eqref{eq:boxed-root-reconstruction}.
The derivative identity says exactly that
\[
 d=1-sQ+\pi s^2.
\]
A calculation in the ring $A$ now gives
\[
 1+xy=t,
 \qquad
 t^2z+\frac{g_1}{g_3}y^2(1+3t)=q,
 \qquad tq=\pi.
\]
The first identity in \eqref{eq:marked-line-identities} gives $B=b$, and the
root equation gives $C=c$.  We have reconstructed a point of the original
fiber.

The two procedures undo one another.  Starting with a source point, one
recovers
\[
 D=t^{-1},\qquad x=S/D,\qquad y=Q-\pi S,\qquad q=\pi D,
\]
and then $z$ from the definition of $q$.  Starting with a root, the marked
coordinate is again $s$.  All formulas use ring operations and the inverse of
the explicitly specified unit $E'(s)$, so they commute with homomorphisms of
test algebras.

The universal property of $K[S]/(E)$ identifies its maps to $A$ with roots of
$E$ in $A$.  Localizing at $\overline{E'}$ imposes exactly the additional
condition that $E'(s)$ be a unit.  This proves
\eqref{eq:localized-natural-fiber-equivalence}, and Yoneda's lemma gives
\eqref{eq:localized-fiber}.  Notice that this general localized statement
requires no separability hypothesis: it removes precisely the ramified part
of the inverse equation.

It remains to check the generic degree.  Put
\[
 L=K(\Pi,B),\qquad \lambda=\frac{g_1}{2},
\]
and write
\[
 E_{\Pi,B,C}(S)=H_{\Pi,B}(S)-\lambda C\in L[C][S].
\]
This polynomial has degree $N$ in $S$.  Its leading coefficient is $g_3\Pi$
when $N=3$ and $g_N\Pi^N$ when $N\geq4$, so it is a unit of $L[C]$.  Hence
$E_{\Pi,B,C}$ is primitive over the UFD $L[C]$.

Assume it factors nontrivially over $L(C)[S]$.  Gauss's lemma then gives a
factorization in $L[C][S]$.  Since the polynomial is linear in $C$, one factor,
say $A(S)$, is independent of $C$.  Comparing the coefficient of $C$ shows
that $A(S)$ divides the nonzero scalar $-\lambda$ in $L[S]$.  Thus $A$ is a
unit, a contradiction.  Therefore $E_{\Pi,B,C}$ is irreducible of degree $N$
over $K(\Pi,B,C)$.

We finish by comparing that root extension with the actual source field.
On the dense chart \(t\neq0\), the change of variables is birational:
\[
 K(x,y,z)=K(\Pi,S,Q).
\]
The forward inclusion follows from the definitions.  For the reverse
inclusion put \(D=1-SQ+\Pi S^2\).  Then
\begin{equation}
 t=D^{-1},\qquad x=\frac{S}{D},\qquad
 y=Q-\Pi S,\qquad q=\Pi D,
 \label{eq:function-field-reconstruction}
\end{equation}
and the defining equation for \(q\) recovers \(z\) rationally.  Since
\[
 B=Q+\beta(\Pi,S),
\]
we have \(K(\Pi,S,Q)=K(\Pi,B,S)\).  The second marked-line identity gives
\[
 C=\frac{2G_\Pi(S)}{g_1}-BS^2.
\]
Because \(K(\Pi,B,S)=K(x,y,z)\) has transcendence degree three and is
generated by the three elements \(\Pi,B,S\), those elements are algebraically
independent.  In particular, over \(L_0=K(\Pi,B)\), the element \(S\) is
transcendental and the displayed expression for \(C\) is a nonconstant
degree-\(N\) polynomial in \(S\).  Hence \(C\) is transcendental over \(L_0\):
otherwise \(S\), which is algebraic over \(L_0(C)\), would be algebraic over
\(L_0\).  Thus \(\Pi,B,C\) are algebraically independent and
\(T=K(\Pi,B,C)\) is the actual target function field.

The evaluation homomorphism from the abstract target rational function field
to the source field is therefore injective.  Under it,
\(E_{\Pi,B,C}\) has the root \(S\), and the reconstruction above gives
\[
 K(x,y,z)=K(\Pi,B,S)=T(S).
\]
The irreducibility just proved makes \(E_{\Pi,B,C}\) the minimal polynomial
of \(S\) over \(T\).  Therefore
\[
 [K(x,y,z):K(\Pi,B,C)]=[T(S):T]=N,
\]
which is exactly \(\gdeg(F_G)=N\).
\end{proof}

\begin{corollary}[Finite-\'etale (squarefree) quadratic-gauge fibers]
\label{cor:squarefree-gauge}
In the setting of Theorem~\ref{thm:localized-gauge}, suppose that
$E=E_{\pi,b,c}$ is squarefree of degree $N$.  Then the localization
disappears and, for every commutative $K$-algebra $A$, there is a natural
bijection
\begin{equation}
 \operatorname{Hom}_{K\text{-alg}}(K[S]/(E),A)
 \simeq \mathcal F_{\pi,b,c}(A).
 \label{eq:natural-fiber-equivalence}
\end{equation}
Consequently,
\[
 F_G^{-1}(\pi,b,c)\simeq\Spec K[S]/(E),
\]
and this is a full Keller fiber of rank \(N\) (in particular, finite
\'etale).
\end{corollary}

\begin{proof}
Since $E$ is squarefree, choose $U,V\in K[S]$ with $UE+VE'=1$.  Modulo $E$,
the class of $V$ is an inverse for $\overline{E'}$, so the localization in
Theorem~\ref{thm:localized-gauge} is just $K[S]/(E)$.  The natural equivalence
and scheme isomorphism follow from that theorem.  Its geometric-degree
statement makes the rank-$N$ fiber full.
\end{proof}

\begin{corollary}[Jacobian-one effective form]\label{cor:effective-gauge}
Let
\[
 L(\Pi,B,C)=\left(\Pi,-\frac{B}{2},C\right),
 \qquad \widetilde F_G=L\circ F_G.
\]
Then
\[
 \det D\widetilde F_G=1,
 \qquad \gdeg(\widetilde F_G)=N.
\]
For every $(\pi,b,c)\in K^*\times K^2$,
\[
 \widetilde F_G^{-1}\left(\pi,-\frac b2,c\right)
 \simeq
 \operatorname{Spec}\left(
   (K[S]/(E_{\pi,b,c}))[1/\overline{E_{\pi,b,c}'}]
 \right).
\]
If $E_{\pi,b,c}$ is squarefree of degree $N$, this reduces to
\[
 \widetilde F_G^{-1}\left(\pi,-\frac b2,c\right)
 \simeq\Spec K[S]/(E_{\pi,b,c}).
\]
In particular, a target with $b=0$ is unchanged by the normalization.
Moreover,
\[
 \deg\Pi\leq7,
 \qquad \deg B\leq6N+2,
 \qquad \deg C\leq6N,
\]
so
\[
 \boxed{\max_i\deg(\widetilde F_G)_i\leq6N+2.}
\]
\end{corollary}

\begin{proof}
The linear map $L$ has determinant $-1/2$, so the determinant, geometric
degree, and localized-fiber statements follow from
Theorem~\ref{thm:localized-gauge}; the squarefree simplification follows from
Corollary~\ref{cor:squarefree-gauge}.

For the degree estimate, $\deg t=2$ and $\deg q\leq5$, hence
$\deg\Pi\leq7$.  A summand of $B$ indexed by $k\geq4$ has degree at most
\[
 2\deg t+(k-2)+k\deg q\leq4+(k-2)+5k=6k+2.
\]
The preceding terms have degrees at most $1$, $6$, and $7$, so
$\deg B\leq6N+2$.  Similarly, $(xq)^k$ has degree at most $6k$, while the
remaining terms of $C$ have degree at most $4$.  Thus $\deg C\leq6N$.
A linear output normalization does not change these bounds.
\end{proof}

\section{Which finite \'etale algebras occur?}

Let $K$ be a field of characteristic zero.  A polynomial map
\[
 F:\A^m_K\longrightarrow\A^m_K
\]
is a \emph{Keller map} if $\det DF\in K^*$.  Its geometric degree is
\[
 \gdeg(F)=[K(x_1,\ldots,x_m):K(F_1,\ldots,F_m)].
\]
A Keller map is \'etale and hence quasi-finite.  Every fiber over a $K$-point
is therefore a finite \'etale $K$-scheme, possibly empty.

\begin{definition}[Keller fiber]\label{def:keller-fiber}
A nonzero finite \'etale $K$-algebra $A$ is a \emph{Keller fiber} if there are
a Keller map $F:\A^m_K\to\A^m_K$ and a point $y\in\A^m(K)$ such that
\[
 F^{-1}(y)\simeq\Spec A,
 \qquad \dim_KA=\gdeg(F).
\]
The equality of the two degrees is the fullness condition.  It says that the
special fiber contains all generic inverse sheets.  Without it, a fiber may be
shorter because some sheets have escaped through the nonproperness boundary.
\end{definition}

The construction used in the proof can be read from left to right:
\[
 \boxed{
 \text{squarefree }P(T)
 \ \longmapsto\
 G(S)=P(a+S)-P(a)
 \ \longmapsto\
 \widetilde F_G^{-1}(y_{P,a})\simeq\Spec K[T]/(P).}
\]
The last arrow is an isomorphism of schemes.  It is not merely a matching of
$K$-points or a count after passing to an algebraic closure.

\begin{lemma}[Degree two is impossible, over every characteristic-zero field]
\label{lem:degree-two}
Let $K$ be a field of characteristic zero, let $m\geq1$, and let
\[
 F=(F_1,\ldots,F_m)\in K[X_1,\ldots,X_m]^m
 \quad\text{satisfy}\quad
 \det\left(\frac{\partial F_i}{\partial X_j}\right)\in K^*.
\]
Then the induced function-field extension cannot have degree two:
\[
 \left[
 K(X_1,\ldots,X_m):K(F_1,\ldots,F_m)
 \right]\ne2.
\]
\end{lemma}

\begin{proof}
We give the scalar-extension and descent steps because the external theorem
used here is most conveniently cited over $\mathbb C$.

\smallskip
\noindent\emph{The precise external input.}
Campbell's unnumbered theorem
\cite[Theorem, p.~244]{Campbell} states that a polynomial map
$H:\mathbb C^m\to\mathbb C^m$ has a polynomial inverse if its Jacobian
determinant never vanishes and the induced extension
\[
 \mathbb C(H_1,\ldots,H_m)\subset
 \mathbb C(X_1,\ldots,X_m)
\]
is normal.  In characteristic zero ``normal'' here is equivalent to
``Galois,'' because the finite function-field extension is separable.
Razar and, independently, Wright later gave algebraic treatments of this
Galois case \cite{Razar,Wright}; the proof below needs only Campbell's stated
complex theorem.

\smallskip
\noindent\emph{Generic degree survives scalar extension.}
Assume for contradiction that the displayed degree over $K$ is two.  Let
$K_0\subset K$ be the subfield generated over $\mathbb Q$ by the finitely
many coefficients of the $F_i$, and write $F_0$ for the same formulas over
$K_0$.  The nonzero Jacobian determinant makes
\begin{equation}
 \varphi_0:K_0[Y_1,\ldots,Y_m]\longrightarrow
 K_0[X_1,\ldots,X_m],
 \qquad Y_i\longmapsto F_i,
 \label{eq:degree-two-coordinate-map}
\end{equation}
injective and generically finite.  Indeed, the Jacobian criterion in
characteristic zero makes the $F_i$ algebraically independent, so both
fraction fields in \eqref{eq:degree-two-coordinate-map} have transcendence
degree $m$ and their extension is finite.  If its degree is $d_0$, there is a
nonempty principal open $U_0$ of the target on which $F_0$ is finite locally
free of rank $d_0$: first clear denominators in the finite extension of
fraction fields, and then apply generic freeness.  For every field extension
$E/K_0$, the base change over the nonempty open $(U_0)_E$ is again finite
locally free of rank $d_0$.  Taking generic fibers gives
\begin{equation}
 [E(X_1,\ldots,X_m):E(F_1,\ldots,F_m)]=d_0.
 \label{eq:degree-two-base-change}
\end{equation}
Applying \eqref{eq:degree-two-base-change} first to $E=K$ shows $d_0=2$.

\smallskip
\noindent\emph{Why degree two gives Campbell's field hypothesis.}
Because $K_0$ is finitely generated over $\mathbb Q$, choose an embedding
$K_0\hookrightarrow\mathbb C$.  Equation
\eqref{eq:degree-two-base-change} with $E=\mathbb C$ says that
\[
 \mathbb C(X_1,\ldots,X_m)/
 \mathbb C(F_1,\ldots,F_m)
\]
has degree two.  It is separable in characteristic zero.  Every separable
quadratic field extension is normal (the second root of the minimal
polynomial of an element is its trace minus that element), hence it is
Galois.  Thus the complexified map satisfies exactly the hypotheses of
Campbell's theorem and is a polynomial automorphism.

\smallskip
\noindent\emph{Faithfully flat descent of the inverse.}
The base change of \eqref{eq:degree-two-coordinate-map} along
$K_0\hookrightarrow\mathbb C$ is therefore an isomorphism.  Since
$\mathbb C$ is faithfully flat over $K_0$, tensoring the kernel and cokernel
of $\varphi_0$ with $\mathbb C$ detects their vanishing.  Hence $\varphi_0$
itself is an isomorphism.  After base change to $K$, the original $F$ is a
polynomial automorphism, so
$K(F_1,\ldots,F_m)=K(X_1,\ldots,X_m)$ and its function-field degree is one.
This contradicts the assumed degree two.
\end{proof}

\begin{theorem}[Prescribed finite-\'etale realization]
\label{thm:main}
Let $P\in K[T]$ be squarefree of degree $N\geq3$.  Choose $a\in K$ outside
the finite zero set of $P'P'''$, and put
\[
 G(S)=P(a+S)-P(a).
\]
Let $F_G=(\Pi,B,C)$ be the quadratic-gauge map of
\eqref{eq:gauge-map}, and define
\[
 \widetilde F_G=L\circ F_G,
 \qquad
 L(\Pi,B,C)=\left(\Pi,-\frac{B}{2},C\right).
\]
Then
\[
 \det D\widetilde F_G=1,
 \qquad \gdeg(\widetilde F_G)=N,
 \qquad \max_i\deg (\widetilde F_G)_i\leq 6N+2.
\]
The Galois closure of the induced degree-$N$ function-field extension has
geometric and arithmetic Galois group $S_N$.
At
\begin{equation}
 y_{P,a}=\left(1,0,-\frac{2P(a)}{P'(a)}\right)
 \label{eq:main-target}
\end{equation}
its full fiber is
\[
 \boxed{\widetilde F_G^{-1}(y_{P,a})
 \simeq\Spec K[T]/(P).}
\]
More explicitly, for every commutative $K$-algebra $R$ there is a bijection
\[
 \operatorname{Hom}_{K\text{-alg}}(K[T]/(P),R)
 \simeq
 \{u\in R^3:\widetilde F_G(u)=y_{P,a}\},
\]
natural in $R$.

Consequently, every finite \'etale $K$-algebra of rank at least three occurs
as a full fiber of a Jacobian-one map in affine three-space.
\end{theorem}

\begin{proof}
The polynomials $P'$ and $P'''$ are nonzero in characteristic zero when
$N\geq3$.  Their product has only finitely many roots, while $K$ is infinite,
so an admissible $a$ exists.  The linear and cubic Taylor coefficients of $G$
are $P'(a)$ and $P'''(a)/3!$, and the quadratic gauge of
Section~\ref{sec:gauge} is therefore defined.

At the target \eqref{eq:main-target}, its inverse polynomial is
\[
 G(S)-\frac{P'(a)}2\left(-\frac{2P(a)}{P'(a)}\right)
 =P(a+S).
\]
Translation preserves squarefreeness, so
Corollary~\ref{cor:squarefree-gauge}
identifies the full fiber with $\Spec K[S]/(P(a+S))$.

The translation back to the original polynomial is explicit:
\[
 K[S]/(P(a+S))\longrightarrow K[T]/(P(T)),
 \qquad \overline S\longmapsto\overline T-a,
\]
with inverse $\overline T\mapsto\overline S+a$.  Composing this with the
natural equivalence of Corollary~\ref{cor:squarefree-gauge} gives, for every
commutative $K$-algebra $R$,
\[
 \operatorname{Hom}_{K\text{-alg}}(K[T]/(P),R)
 \simeq \widetilde F_G^{-1}(y_{P,a})(R).
\]
Thus the isomorphism includes the complete scheme structure and commutes with
base change.

The linear map $L$ fixes the chosen target because its second coordinate is
zero.  It multiplies the Jacobian by $-1/2$ and changes neither geometric
degree nor coordinate degree.  Corollary~\ref{cor:effective-gauge} therefore
gives the determinant-one statement and the bound $6N+2$.  Since both the
quotient algebra and the generic inverse problem have degree $N$, the fiber is
full.  The symmetric-monodromy assertion follows from
Theorem~\ref{thm:symmetric-monodromy}.

For a general finite \'etale algebra of rank at least three, use the
monogenicity lemma in Section~\ref{sec:realization} and apply the polynomial
case just proved.
\end{proof}

\begin{theorem}[Classification of full Keller fibers]
\label{cor:rank-classification}
For $N\geq1$, let $\mathsf{FullFib}_{3,N}(K)$ be the replete full subcategory
of finite \'etale $K$-schemes whose objects are the schemes isomorphic to
full fibers
\[
 F^{-1}(y),\qquad
 F:\A^3_K\longrightarrow\A^3_K,\quad
 \gdeg(F)=N,\quad y\in\A^3(K).
\]
Then
\[
 \mathsf{FullFib}_{3,N}(K)=
 \begin{cases}
  \mathsf{FinEt}^{\,N}_K,&N=1\text{ or }N\geq3,\\
  \varnothing,&N=2,
 \end{cases}
\]
where $\mathsf{FinEt}^{\,N}_K$ is the full subcategory of rank-$N$ finite
\'etale $K$-schemes.  In particular, for $N\geq3$, a finite $K$-scheme occurs
as a full fiber of a geometric-degree-$N$ Keller map in $\A^3_K$ if and only
if it is finite \'etale of rank $N$.

Equivalently on coordinate algebras, every rank-$N$ finite \'etale
$K$-algebra occurs for $N=1$ or $N\geq3$, and none occurs for $N=2$.
Consequently, the possible nonzero ranks of full Keller fibers are exactly
\[
 \boxed{1,3,4,5,\ldots}.
\]
\end{theorem}

\begin{proof}
The Jacobian condition makes $F$ \'etale.  Hence every fiber over a
$K$-point is a finite \'etale $K$-scheme, and the fullness equality makes its
rank equal to $\gdeg(F)=N$.  This proves the necessary direction in every
degree.  The identity map realizes the unique rank-one object.
Theorem~\ref{thm:main} realizes every rank-$N$ finite \'etale object in
affine three-space for every $N\geq3$, while
Lemma~\ref{lem:degree-two} shows that no geometric-degree-two Keller map
exists.  This last exclusion concerns the generic function-field degree
before any target point is selected.  Thus an empty special fiber cannot
create a rank-two object; it is not a nonzero finite \'etale scheme at all.
Nor can a shorter special fiber of a map of some other geometric degree:
such a fiber fails the defining fullness equality and is not an object of
$\mathsf{FullFib}_{3,2}(K)$.  Empty and non-full fibers therefore have no
bearing on the asserted rank set.
\end{proof}

\begin{remark}[What the categorical statement remembers]
The category $\mathsf{FullFib}_{3,N}(K)$ is the essential image inside finite
\'etale $K$-schemes: it remembers the resulting fiber scheme, not a chosen
presentation by a map--target pair $(F,y)$.  Different pairs may realize the
same object.  If one works with finite \'etale algebras rather than schemes,
the corresponding statement is contravariant under $\Spec$.
\end{remark}


\section{Symmetric monodromy}
\label{sec:monodromy}

\begin{lemma}[Morse slices of the quadratic pencil]
\label{lem:quadratic-pencil-morse}
Let \(k\) be an algebraically closed field of characteristic zero, and let
\[
 H(S)=h_1S+\cdots+h_NS^N,\qquad h_1h_N\neq0.
\]
Then
\[
 f_b(S)=H(S)-\frac{h_1b}{2}S^2
\]
is Morse for all but finitely many \(b\in k\).
\end{lemma}

\begin{proof}
A repeated root \(r\) of
\[
 H(S)-\frac{h_1}{2}(BS^2+C)
\]
cannot be zero, because \(H'(0)=h_1\).  Solving \(E(r)=E'(r)=0\) gives
\begin{equation}
 B(r)=\frac{H'(r)}{h_1r},
 \qquad
 C(r)=\frac{2H(r)-rH'(r)}{h_1},
 \qquad r\in\mathbb G_m.
 \label{eq:quadratic-discriminant-parametrization}
\end{equation}
Their differentials satisfy
\begin{equation}
 dC=-r^2\,dB.
 \label{eq:quadratic-legendre}
\end{equation}
Let \(M=k(B(r),C(r))\subset k(r)\).  Equation
\eqref{eq:quadratic-legendre}, with \(B\) nonconstant, shows that
\(r^2\in M\).  Since
\([k(r):k(r^2)]=2\), either \(M=k(r^2)\) or \(M=k(r)\).  The first
possibility is excluded by the pole of order one of \(B(r)\) at \(r=0\):
every element of \(k(r^2)\) has even valuation there.  Thus \(M=k(r)\), so
\eqref{eq:quadratic-discriminant-parametrization} is birational onto its
image.

The rational map \(r\mapsto B(r)\) has degree \(N-1\).  Its only poles are
the pole of order one at \(0\) and the pole of order \(N-2\) at infinity;
the latter is present because \(h_N\neq0\).  Outside finitely many values of
\(b\), the fiber \(B(r)=b\) therefore consists of \(N-1\) distinct points.
They are precisely the critical points of \(f_b\).

At such a point,
\[
 f_b(r)=H(r)-\frac{rH'(r)}2=\frac{h_1}{2}C(r).
\]
The normalization map in
\eqref{eq:quadratic-discriminant-parametrization} is an isomorphism away
from finitely many points of the image curve.  After excluding the finitely
many \(B\)-coordinates of its exceptional points, no two points in a fiber
of \(B\) have the same \(C\)-coordinate.  Thus the critical values of
\(f_b\) are pairwise distinct, and \(f_b\) is Morse.
\end{proof}

\begin{lemma}[Specialization through the ordered-root cover]
\label{lem:slice-detects-monodromy}
Let \(k\) be a field, let \(U\) be a connected normal \(k\)-variety, and let
\[
 f(S)\in\Gamma(U,\mathcal O_U)[S]
\]
have degree \(N\), with leading coefficient and discriminant invertible on
\(U\).  Let \(G\leq S_N\) be the Galois group of \(f\) over \(k(U)\), acting
on a chosen ordering of its roots.  If \(Z\) is a connected normal
\(k\)-variety and \(i:Z\to U\) is a morphism, then the Galois group of
\(i^*f\) over \(k(Z)\), in its root action, is conjugate in \(S_N\) to a
subgroup of \(G\).

In particular, let \(k\) be algebraically closed, fix \(p\in k^*\) and
\(b\in k\), and suppose that \(E_{p,b,C}(S)\) is separable over \(k(C)\).
Then its splitting-field Galois group is conjugate to a subgroup of the
Galois group of \(E_{\Pi,B,C}(S)\) over \(k(\Pi,B,C)\).
\end{lemma}

\begin{proof}
After dividing by the invertible leading coefficient, form the finite
\'etale \(U\)-scheme of ordered \(N\)-tuples of pairwise distinct roots.
Choose the connected component \(Y\) containing the chosen generic
ordering.  The induced cover \(Y\to U\) is the splitting cover: it is a
finite \'etale \(G\)-torsor, and its function field is the splitting field
of \(f\).

The pullback \(Y_Z=Y\times_U Z\) is again a finite \'etale \(G\)-torsor,
although it need not be connected.  Choose a connected component \(Y_0\).
The subgroup
\[
 H=\{g\in G:gY_0=Y_0\}
\]
acts simply transitively on the geometric fibers of \(Y_0\to Z\).
Consequently \(k(Y_0)/k(Z)\) is the splitting field of \(i^*f\) and
\(\operatorname{Gal}(k(Y_0)/k(Z))=H\leq G\).  A different ordering or a
different component conjugates \(H\) in \(S_N\), which proves the first
statement.

For the asserted specialization, take
\[
 U=\Spec k[\Pi,\Pi^{-1},B,C,\Delta_G^{-1}]
\]
and divide \(E_{\Pi,B,C}\) by its invertible leading coefficient.  Since
\(\Delta_G(p,b,C)\) is a nonzero polynomial, the open curve
\[
 Z=\Spec k[C,\Delta_G(p,b,C)^{-1}]
\]
is connected and maps to \(U\) by \(\Pi\mapsto p,\ B\mapsto b\).  The first
statement gives the required subgroup inclusion.
\end{proof}

\begin{theorem}[Full symmetric generic monodromy]
\label{thm:symmetric-monodromy}
For every seed \eqref{eq:seed}, not merely for a generic choice of its
coefficients, the geometric and arithmetic monodromy groups of $F_G$ are
\[
 \boxed{\operatorname{Mon}_{\mathrm{geom}}(F_G)
       =\operatorname{Mon}_{\mathrm{arith}}(F_G)=S_N.}
\]
Equivalently, the splitting field of $E_{\Pi,B,C}(S)$ over
$K(\Pi,B,C)$ has Galois group $S_N$, and the same is true after replacing
$K$ by an algebraic closure.  The conclusion is unchanged for the
Jacobian-one normalization $\widetilde F_G$.
\end{theorem}

\begin{proof}
We work first over $\overline K$.  Fix $p\in\overline K^*$ and put
$H(S)=G_p(S)$.  Consider
\[
 f_b(S)=H(S)-\frac{g_1b}{2}S^2.
\]
Lemma~\ref{lem:quadratic-pencil-morse} shows that \(f_b\) is Morse for all
but finitely many \(b\in\overline K\).

Choose one such $b$.  The classical Morse-polynomial theorem
\cite[Theorem~4.4.5]{Serre} gives
\[
 \operatorname{Gal}\left(
 f_b(S)-T\,/\,\overline K(T)\right)=S_N.
\]
Indeed, the polynomial cover is connected, its finite branch inertia groups
are transpositions, and those inertia groups generate; a transitive
permutation group generated by transpositions is the full symmetric group.
After the substitution $T=g_1C/2$, this is the inverse cover on the line
\(\Pi=p,\ B=b\).  Lemma~\ref{lem:slice-detects-monodromy} embeds its
monodromy into the generic geometric monodromy.  Hence
\[
 S_N\leq\operatorname{Mon}_{\mathrm{geom}}(F_G)\leq S_N,
\]
which proves equality.  Finally,
\[
 \operatorname{Mon}_{\mathrm{geom}}(F_G)
 \triangleleft\operatorname{Mon}_{\mathrm{arith}}(F_G)\leq S_N
\]
forces the arithmetic group to be $S_N$ as well.  A linear automorphism of
the target does not change either monodromy group, so the same conclusion
holds for $\widetilde F_G$.
\end{proof}

\begin{corollary}[Absolute and stable compositional atomicity]
\label{cor:stable-atomicity}
Let \(F\) be any realization map \(F_G\) or \(\widetilde F_G\) above.  If
\[
 F=H_2\circ H_1
\]
as polynomial self-maps of \(\A^3_K\), then one factor is a polynomial
automorphism.  The same conclusion holds after every characteristic-zero
extension of \(K\), after every stabilization
\[
 F\times\operatorname{id}_{\A^r},
\]
and after polynomial changes of source and target coordinates.  Thus every
realization map is absolutely and stably compositionally indecomposable.
\end{corollary}

\begin{proof}
Over an algebraic closure, write
\[
 K_0=\overline K(F_1,F_2,F_3)
 \subset L=\overline K(x,y,z).
\]
The generic sheets form the natural \(S_N\)-set.  This action is primitive:
its point stabilizer \(S_{N-1}\) is maximal in \(S_N\).  Galois
correspondence therefore gives no proper intermediate field between \(K_0\)
and \(L\).

Any factorization \(F=H_2\circ H_1\) gives such an intermediate field.
Both factors are dominant, and the chain rule
\[
 \det DF=(\det DH_2)\circ H_1\cdot\det DH_1\in K^*
\]
shows that both are Keller maps: the only units in a polynomial ring are
constants, and dominance of \(H_1\) makes pullback injective.  Primitivity
forces one factor to have geometric degree one.  A degree-one Keller
self-map of affine space is an automorphism: it is birational and \'etale,
hence an open immersion by Zariski's Main Theorem; a proper complement is
excluded first in codimension one by the unit group of affine space and then
in codimension at least two by normal Hartogs extension.

Extension of the constant field does not change the geometric monodromy
action.  Stabilization replaces \(K_0\subset L\) by
\[
 K_0(t_1,\ldots,t_r)\subset L(t_1,\ldots,t_r),
\]
with the same Galois closure, group, and point stabilizer.  The preceding
argument therefore applies after every extension and stabilization.
Conjugating by polynomial automorphisms plainly preserves the conclusion.
\end{proof}

\begin{remark}
Taking \(N=12\) gives an atomic Keller map of composite geometric degree.
An explicit sparse formula and independent certificate are recorded
\href{https://github.com/royvanrijn/jacobian-research/blob/main/verified/INDECOMPOSABLE_COMPOSITE_DEGREE.md}
{separately in the accompanying repository}.
\end{remark}

\begin{corollary}[Hilbertian specializations]\label{cor:hilbert-monodromy}
Suppose that $K$ is Hilbertian.  For every seed \eqref{eq:seed}, there are
infinitely many targets in $K^*\times K^2$ whose fibers are connected full
finite-\'etale schemes of rank $N$ and whose splitting fields have Galois
group $S_N$ over $K$.
\end{corollary}

\begin{proof}
Choose \(p\in K^*\).  Lemma~\ref{lem:quadratic-pencil-morse} excludes only
finitely many values of \(b\) in an algebraic closure.  A Hilbertian field is
infinite, so some \(b\in K\) makes \(f_b\) Morse.  Let \(L\) be the splitting
field of \(f_b(S)-T\) over \(K(T)\).  Serre's Morse-polynomial theorem gives
\[
 \operatorname{Gal}(L/K(T))=S_N
\]
\cite[Theorem~4.4.5]{Serre}.  Applied after extending the constants to
\(\overline K\), the same theorem gives geometric group \(S_N\).  Thus the
arithmetic and geometric groups agree, their constant-field quotient is
trivial, and \(L/K(T)\) is regular over \(K\).

Hilbert irreducibility in its Galois form
\cite[Chapter~3]{Serre} now supplies infinitely many $c\in K$ for which
$E_{p,b,c}$ remains irreducible and has Galois group $S_N$.  Such a polynomial
is squarefree of degree $N$, so Corollary~\ref{cor:squarefree-gauge}
identifies it with a connected full Keller fiber.
\end{proof}

\begin{remark}[Which hypotheses are used where]\label{rem:gauge-hypotheses}
The general localized quadratic-gauge fiber theorem,
Theorem~\ref{thm:localized-gauge}, uses only that $\pi$, $g_1$, and $g_3$ are
nonzero for the fiber identification, together with the identity
\[
 E'(S)=g_1\bigl(1-bS+S\beta(\pi,S)+\pi S^2\bigr).
\]
Squarefreeness is needed only to conclude that $E'$ is invertible at every root
and hence that no localization remains.  The condition $g_3\neq0$ is what
makes the source coefficient $g_1/g_3$ meaningful.  For a translated
polynomial $P(a+S)-P(a)$, the cubic coefficient is the third Hasse derivative
at $a$, equal to $P'''(a)/3!$ in characteristic zero.

The transposition inertia in Theorem~\ref{thm:symmetric-monodromy} does not
contradict the fact that $F_G$ is \'etale.  A repeated inverse root has
$E'=0$, so the reconstruction theorem removes it from the affine source;
the corresponding valuation is centered on the nonproperness boundary.
Monodromy is taken on the function-field cover, or equivalently on the finite
\'etale restriction over a sufficiently small target open set.
\end{remark}
